\documentclass[10pt,a4paper]{amsart}
\usepackage[margin=19mm]{geometry}
\usepackage{amsmath,amssymb,mathtools}
\usepackage[scr=rsfs,cal=euler]{mathalfa}
\usepackage{xfrac}
\usepackage[unicode,hidelinks,bookmarks=false]{hyperref}
\newcommand{\ie}{i.e.\ }
\newcommand{\cf}{cf.\ }
\newcommand{\resp}{respectively\ }
\newcommand{\Subsection}{\subsection}
\providecommand{\nobreakdash}{\nobreak\hbox{-}}
\setlength{\emergencystretch}{2em}
\multlinegap0pt
\usepackage{amssymb}
\newtheorem{prop}{Proposition}
\newtheorem{thm}{Theorem}
\newcommand{\C}{\mathbb{C}}
\def\D{\mathbb{D}}
 \DeclareMathOperator{\supp}{supp}
\title{Indeterminate Stieltjes moment problems revisited}
\author{Christian Berg}
\date{Source: arXiv:2407.08511v1 (CC BY 4.0)}
\begin{document}
\maketitle

\noindent\emph{Source and licence.} Indeterminate Stieltjes moment problems revisited. Version arXiv:2407.08511v1 (CC BY 4.0), distributed by arXiv under the Creative Commons Attribution 4.0 International licence, http://creativecommons.org/licenses/by/4.0/, as recorded in the arXiv licence metadata for this exact version and shown on the abstract page https://arxiv.org/abs/2407.08511v1. Full licence notice: \texttt{licenses/arxiv-2407.08511-CC-BY-4.0.txt}.\par\smallskip

\noindent\textit{Editorial scope.} Counted unit: the article's thm thm:Fr2 (source0.tex lines 424--426). The article's complete original proof of it is delivered with it (source0.tex lines 428--456, one \texttt{proof} environment of 29 lines). No mathematics is written, repaired or re-derived: the statement and the proof are the article's own lines, in the article's own notation, and no retained line is substituted or re-flowed.  Retained uncounted as the source-local context the proof needs: lines 245--256; lines 418--420.  Nothing was omitted from any retained span, so the package carries no omission record.  The retained lines cite 1 works, whose entries are transcribed verbatim from the article's own bibliography source in the e-print and delivered with short editorial labels recorded in the item's reference mapping.  No retained line contains a figure environment, an image-inclusion command or drawing code, so no image is delivered and the package declares no asset dependency.  Editorial formatting only, in addition: an AMS \texttt{amsart} preamble that restates the article's own declarations and macros used by the retained lines, namely the environments \texttt{prop}, \texttt{thm} on separate counters, together with the standard packages those lines need.  The retained environments print in this document as Proposition 1, Theorem 1 in order of appearance, whereas the article prints the same items with the numbering of its own full document.  The complete native source of the article as distributed in the arXiv e-print for this exact version is bundled unchanged as \texttt{sources/162182/source0.tex}.
\begin{prop}\label{thm:2to1} For $u, v\in\C$ we have
 \begin{eqnarray*}
A_n(u,v)&=& \left|\begin{array}{cc}
 A_n(u)&\;A_n(v)\\C_n(u)&\;C_n(v)\end{array}\right|\\
 B_n(u,v)&=& \left|\begin{array}{cc}
 B_n(u)&\;A_n(v)\\D_n(u)&\;C_n(v)\end{array}\right|\\
  C_n(u,v)&=& \left|\begin{array}{cc}
 A_n(u)&\;B_n(v)\\C_n(u)&\;D_n(v)\end{array}\right|\\
  D_n(u,v)&=& \left|\begin{array}{cc}
 B_n(u)&\;B_n(v)\\D_n(u)&\;D_n(v)\end{array}\right|.
 \end{eqnarray*}
\end{prop}

\begin{equation}\label{eq:el1}
\xi_1=\xi(\nu_{\alpha(s)}),\quad \alpha(s)=D(\xi_1)/B(\xi_1).
\end{equation}

\begin{thm}\label{thm:Fr2} Let $s$ be a normalized Stieltjes moment sequence which is indet(S). For any solution $\nu\in V$ to the indeterminate Hamburger problem we have
$\xi(\nu)<\xi(\nu_{\alpha(s)})$ unless $\nu$ is the Friedrichs solution $\nu_{\alpha(s)}$.
\end{thm}

\begin{proof} Suppose that $\nu$ is a solution satisfying $\supp(\nu)\subseteq [\xi_1,\infty)$, and we want to show that $\nu=\nu_{\alpha(s)}$.

Considering the translation $\tau(x)=x-\xi_1$,  the image measures $\tau(\nu_{\alpha(s)}), \tau(\nu)$ are supported on $[0,\infty)$ and belong to the Stieltjes moment problem  with moments
$$
\tilde{s}_n=\int (x-\xi_1)^n\,d\nu(x)=\sum_{j=0}^n\binom{n}{j} (-\xi_1)^js_{n-j}.
$$
It is clearly indeterminate as Hamburger moment problem and the full set of solutions is
$\{\tau(\mu)\mid \mu\in V\}$. The orthonormal polynomials and those of the second kind are $P_n(x+\xi_1),Q_n(x+\xi_1)$. The corresponding Nevanlinna functions marked with  $\tilde{}$ are
$$
\tilde{A}(u,v)=A(u+\xi_1,v+\xi_1),\ldots,\tilde{D}(u,v)=D(u+\xi_1,v+\xi_1),\quad u,v\in\C
$$

From the formulas for the Nevanlinna functions of two variables expressed by the Nevanlinna functions of one variable, see Proposition~\ref{thm:2to1} (where $n\to\infty$), we get
\begin{eqnarray*}
\tilde{A}(z)&=&A(z+\xi_1)C(\xi_1)-C(z+\xi_1)A(\xi_1)\\
\tilde{B}(z)&=&B(z+\xi_1)C(\xi_1)-D(z+\xi_1)A(\xi_1)\\
\tilde{C}(z)&=&A(z+\xi_1)D(\xi_1)-C(z+\xi_1)B(\xi_1)\\
\tilde{D}(z)&=&B(z+\xi_1)D(\xi_1)-D(z+\xi_1)B(\xi_1).
\end{eqnarray*}
These formulas were found in \cite[Proposition 3.3]{Pe1}.
According to these formulas we get using \eqref{eq:el1}
$$
\frac{\tilde{D}(x)}{\tilde{B}(x)}=\frac{D(\xi_1)-B(\xi_1)\frac{D(x+\xi_1)}{B(x+\xi_1)}}
{C(\xi_1)-A(\xi_1)\frac{D(x+\xi_1)}{B(x+\xi_1)}}\to \frac{D(\xi_1)-B(\xi_1)\alpha(s)}
{C(\xi_1)-A(\xi_1)\alpha(s)}=0
$$
for $x\to -\infty$,  and this shows that $\tilde{\alpha}:=\alpha(\tilde{s})=0$ for the new Stieltjes sequence $\tilde{s}$, so $\tilde{s}$ is det(S). We conclude that $\tau(\nu_{\alpha(s)})=\tau(\nu)$
and hence $\nu_{\alpha(s)}=\nu$.
\end{proof}

\begin{thebibliography}{99}
\bibitem[Pe1]{Pe1} H.~L.~Pedersen, Stieltjes Moment Problems and the Friedrichs Extension of a Positive Definite Operator, J. Approx. Theory {\bf 83} (1995), 289--307.
\end{thebibliography}
\end{document}
